\chapter{Examples and convention tests}
\label{chap:examples}

The examples in this chapter are computational proofs as well as
regression tests.  Each isolates a convention that can otherwise remain
incorrect for many pages: geometric versus arithmetic Frobenius, the
alternating determinant sign, compact supports, the direction of a Tate
twist, geometric connectedness, tensor-product weights, or base-field
Frobenius.  Every answer is computed independently from point counts and
from cohomology.

\section{The rational point}

Let \(X_0=\Spec\Fq\).  There is one point over every finite extension, so
\[
 N_n=1,\qquad
 Z(X_0,T)=\exp\left(\sum_{n\geq1}\frac{T^n}{n}\right)
          =\frac1{1-T}.
 \tag{15.1}\label{eq:zeta-point}
\]
The closed-point calculation has one degree-one point and gives the same
Euler factor.

Cohomologically,
\[
 H^0(X,\Ql)=\Ql,\qquad H^i(X,\Ql)=0\quad(i>0),
\]
and geometric Frobenius acts trivially on constants.  Hence
\[
 P_0(T)=1-T.
\]
Because \(d=0\), degree zero is also top degree; the endpoint factor is
one polynomial, not two copies of it.

The Euler characteristic is \(1\), and the functional equation becomes
\[
 Z(T^{-1})=-T Z(T).
\]
Indeed,
\[
 \frac1{1-T^{-1}}=-\frac{T}{1-T}.
\]
Thus the sign \(\varepsilon\) in the general theorem can be negative even
for the simplest smooth proper variety.

\section{The affine line and compact supports}

For \(X_0=\A^1_{\Fq}\),
\[
 N_n=q^n,\qquad
 Z(\A^1,T)=\frac1{1-qT}.
 \tag{15.2}\label{eq:zeta-A1}
\]
The scheme is smooth but not proper.  Ordinary cohomology sees only
\(H^0=\Ql\) and would produce \(1/(1-T)\), which is wrong.  Compactly
supported cohomology is
\[
 H_c^i(\A^1_{\Fbar},\Ql)=
 \begin{cases}
   \Ql(-1),&i=2,\\
   0,&i\neq2.
 \end{cases}
\]
Geometric Frobenius acts on \(\Ql(-1)\) by \(q\), and the even
cohomological degree appears in the denominator:
\[
 \det(1-TF\mid H_c^2)^{-1}=(1-qT)^{-1}.
\]

\begin{convention}
The calculation \eqref{eq:zeta-A1} tests three choices simultaneously:
\begin{itemize}
\item the trace formula uses geometric Frobenius;
\item \(F\) acts by \(q\) on \(\Ql(-1)\);
\item the exponent in the zeta product is \((-1)^{i+1}\).
\end{itemize}
Reversing any one choice gives either \(1/(1-q^{-1}T)\) or
\(1-qT\), both incompatible with the elementary point count.
\end{convention}

\section{Projective spaces}

The decomposition of \(\PP^r(\mathbf F_{q^n})\) into affine cells gives
\[
 \#\PP^r(\mathbf F_{q^n})=1+q^n+\cdots+q^{rn}.
\]
Taking the exponential term by term,
\[
 Z(\PP^r,T)=\prod_{j=0}^r(1-q^jT)^{-1}.
 \tag{15.3}\label{eq:zeta-projective-space}
\]
Etale cohomology gives the identical answer:
\[
 H^{2j}(\PP^r_{\Fbar},\Ql)=\Ql(-j)
 \quad(0\leq j\leq r),\qquad
 H^{2j+1}=0.
\]
Thus geometric Frobenius has eigenvalue \(q^j\) in degree \(2j\), of
weight \(2j\), and every nonzero group contributes a denominator factor.

For \(r=1\),
\[
 Z(\PP^1,T)=\frac1{(1-T)(1-qT)}.
\]
Here \(d=1\), \(b_0=b_2=1\), and \(\chi=2\).  A direct substitution gives
\[
 Z(\PP^1,q^{-1}T^{-1})=qT^2Z(\PP^1,T).
\]
This tests both endpoint factors and the positive sign for the
projective line.

For general \(r\), the transformation \(j\mapsto r-j\) pairs the factors
\(1-q^jT\) and \(1-q^{r-j}T\), visibly reproducing Poincare duality and
the multiplier \(q^r\).

\section{Finite etale schemes and component permutations}

Every finite etale \(\Fq\)-scheme has a unique presentation
\[
 Y_0=\coprod_{a=1}^s\Spec\mathbf F_{q^{m_a}}.
\]
There is one closed point of degree \(m_a\) for each component, so
\[
 Z(Y_0,T)=\prod_{a=1}^s(1-T^{m_a})^{-1}.
 \tag{15.4}\label{eq:zeta-finite-etale}
\]
Equivalently,
\[
 \#Y_0(\mathbf F_{q^n})
 =\sum_{a:m_a\mid n}m_a,
\]
because a field embedding
\(\mathbf F_{q^{m_a}}\hookrightarrow\mathbf F_{q^n}\) exists precisely
when \(m_a\mid n\), and then there are \(m_a\) of them.

The geometric fiber is a finite set \(S\).  Its cohomology is concentrated
in degree zero:
\[
 H^0(Y,\Ql)=\Ql^S.
\]
On the points arising from \(\mathbf F_{q^m}\), Frobenius is a cycle of
length \(m\).  The eigenvalues of its permutation matrix are all
\(m\)-th roots of unity, once each.  Therefore
\[
 \det(1-TF\mid\Ql^{m})
 =\prod_{\zeta^m=1}(1-\zeta T)=1-T^m.
\]
Multiplication over the cycles recovers
\eqref{eq:zeta-finite-etale}.

This is the cleanest test of geometric connectedness.
\(\Spec\mathbf F_{q^m}\) is connected as an \(\Fq\)-scheme, but its
geometric fiber consists of \(m\) points.  For \(m>1\),
\[
 P_0(T)=1-T^m\neq1-T.
\]
A finite etale scheme over \(\Fq\) is geometrically connected only in
the one-point case \(\Spec\Fq\).

\section{The multiplicative group}

For \(\Gm=\A^1-\{0\}\),
\[
 \#\Gm(\mathbf F_{q^n})=q^n-1,
\qquad
 Z(\Gm,T)=\frac{1-T}{1-qT}.
 \tag{15.5}\label{eq:zeta-Gm}
\]
We now derive both compact-support groups rather than merely quote them.
For the open immersion \(j:\Gm\hookrightarrow\A^1\) and closed
complement \(i:\{0\}\hookrightarrow\A^1\), the localization triangle
\[
 j_!\Ql\longrightarrow\Ql\longrightarrow i_*\Ql\longrightarrow
\]
gives a long exact compact-support sequence.  Since
\[
 H_c^\bullet(\A^1)=\Ql(-1)[-2],
\qquad
 H_c^\bullet(\{0\})=\Ql[0],
\]
the connecting map identifies
\[
 H_c^1(\Gm,\Ql)=\Ql,
\qquad
 H_c^2(\Gm,\Ql)=\Ql(-1),
\]
and all other groups vanish.  The construction is Frobenius equivariant.
Frobenius therefore acts by \(1\) in degree \(1\) and by \(q\) in degree
\(2\).  The alternating determinant is
\[
 \frac{\det(1-TF\mid H_c^1)}
      {\det(1-TF\mid H_c^2)}
 =\frac{1-T}{1-qT},
\]
which is \eqref{eq:zeta-Gm}.  The example proves that an odd
compact-support group can create a numerator factor even for a smooth
one-dimensional variety.

\section{Smooth proper curves}

Let \(C_0/\Fq\) be smooth, proper, geometrically connected, of genus
\(g\).  Its etale cohomology has dimensions
\[
 b_0=1,\qquad b_1=2g,\qquad b_2=1.
\]
Consequently
\[
 Z(C_0,T)=\frac{P_1(T)}{(1-T)(1-qT)},
\quad
 P_1(T)=\det(1-TF\mid H^1(C,\Ql)).
 \tag{15.6}\label{eq:zeta-curve}
\]
The theorem gives \(P_1\in\mathbf Z[T]\), independent of \(\ell\), and
every reciprocal root has absolute value \(q^{1/2}\) under every complex
embedding.

The cup-product pairing on \(H^1\) is alternating and satisfies
\[
 \langle Fx,Fy\rangle=q\langle x,y\rangle.
\]
Thus eigenvalues occur in pairs
\[
 \alpha,\quad q/\alpha.
\]
The determinant of a \(2g\)-dimensional symplectic similitude of
multiplier \(q\) is \(q^g\).  Applying
\cref{prop:reciprocal-characteristic} to the middle pairing gives
\[
 P_1(q^{-1}T^{-1})=q^{-g}T^{-2g}P_1(T).
 \tag{15.7}\label{eq:curve-P1-reciprocal}
\]
Substitution into \eqref{eq:zeta-curve}, with
\(\chi=2-2g\), yields
\[
 Z(C_0,q^{-1}T^{-1})
 =q^{1-g}T^{2-2g}Z(C_0,T).
\]
Thus \(\varepsilon=1\) for every smooth proper geometrically connected
curve.

The trace formula in degree one is
\[
 \#C_0(\Fq)=1+q-\Tr(F\mid H^1).
\]
If the eigenvalues are \(\alpha_1,\ldots,\alpha_{2g}\), then every
complex embedding has \(|\alpha_j|=q^{1/2}\).  The triangle inequality
therefore gives
\[
 \left|\#C_0(\Fq)-(q+1)\right|\leq2gq^{1/2}.
 \tag{15.8}\label{eq:hasse-weil-bound}
\]
For \(g=0\) this recovers the projective-line count.  For \(g=1\), write
\[
 P_1(T)=1-aT+qT^2,
\]
where \(a=q+1-\#C_0(\Fq)\); then
\eqref{eq:hasse-weil-bound} is \(|a|\leq2\sqrt q\).

In this book the curve bound is a consequence of the terminal
smooth-proper purity theorem.  It is not used as the first step of a
Lefschetz-pencil proof.

\section{Products and Kunneth weights}

Let \(X_0\) and \(Y_0\) be smooth and proper.  Pointwise,
\[
 \#(X_0\times Y_0)(\mathbf F_{q^n})
 =\#X_0(\mathbf F_{q^n})\#Y_0(\mathbf F_{q^n}).
 \tag{15.9}\label{eq:product-counts}
\]
This does not make zeta functions multiply: the logarithmic coefficient
of a zeta function is the point count, and
\eqref{eq:product-counts} is a product of those coefficients.

The Kunneth isomorphism is Frobenius equivariant:
\[
 H^r(X\times Y,\Ql)
 \simeq\bigoplus_{i+j=r}
 H^i(X,\Ql)\otimes H^j(Y,\Ql).
 \tag{15.10}\label{eq:kunneth-example}
\]
If \(\alpha\) and \(\beta\) are eigenvalues on the two factors, the
tensor eigenvalue is \(\alpha\beta\).  If their weights are \(i\) and
\(j\), then for every complex realization
\[
 |\iota(\alpha\beta)|
 =q^{i/2}q^{j/2}=q^{(i+j)/2}.
\]
Thus weights add exactly as cohomological degrees add.  Taking
dimensions with alternating signs in \eqref{eq:kunneth-example} gives
\[
 \chi(X_0\times Y_0)=\chi(X_0)\chi(Y_0).
\]
Poincare multipliers multiply as well:
\(q^{\dim X}q^{\dim Y}=q^{\dim X+\dim Y}\).

For example,
\[
 Z(\PP^1\times\PP^1,T)
 =\frac1{(1-T)(1-qT)^2(1-q^2T)}.
\]
The square on \(1-qT\) records the two generators of \(H^2\), one from
each ruling.  This is a useful multiplicity and tensor-product test.

\section{Extension of the finite base field}

Fix \(m\geq1\), and put
\[
 X_m=X_0\times_{\Fq}\mathbf F_{q^m}.
\]
Geometric Frobenius for the new ground field is \(F^m\).  If the
eigenvalues of \(F\) on \(H^i\) are \(\alpha_{ij}\), then
\[
 P_i(X_m,T)
 =\det(1-TF^m\mid H^i(X,\Ql))
 =\prod_j(1-\alpha_{ij}^mT).
 \tag{15.11}\label{eq:base-change-polynomial}
\]
Purity is preserved with the same numerical weight:
\[
 |\iota(\alpha_{ij}^m)|=q^{mi/2}=(q^m)^{i/2}.
\]
The top eigenvalue becomes \(q^{md}\), and the functional equation uses
\((q^m)^{-d}\).  Betti numbers and Euler characteristic do not change.

There is also a point-count identity.  Let \(\mu_m\) be the set of
complex \(m\)-th roots of unity.  Then
\[
 Z(X_m,T^m)=\prod_{\zeta\in\mu_m}Z(X_0,\zeta T).
 \tag{15.12}\label{eq:base-change-zeta}
\]
Taking logarithms, the right side becomes
\[
\begin{aligned}
 \sum_{\zeta^m=1}\sum_{n\geq1}
 \#X_0(\mathbf F_{q^n})\frac{\zeta^nT^n}{n}
 &=\sum_{r\geq1}
 \#X_0(\mathbf F_{q^{mr}})\frac{T^{mr}}r,
\end{aligned}
\]
because the root-of-unity sum is \(m\) when \(m\mid n\) and zero
otherwise.  The last expression is \(\log Z(X_m,T^m)\), proving
\eqref{eq:base-change-zeta}.

\begin{formalization}
These calculations should be executable tests with the following
expected outputs:
\begin{enumerate}
\item `Spec Fq`: \(P_0=1-T\), and the dimension-zero functional-equation
      sign is \(-1\).
\item `A1`: \(H_c^2=\Ql(-1)\), geometric Frobenius has eigenvalue \(q\),
      and \(Z=(1-qT)^{-1}\).
\item `ProjectiveSpace r`: the even-degree eigenvalues are
      \(1,q,\ldots,q^r\), giving \eqref{eq:zeta-projective-space}.
\item finite etale schemes: a cycle of length \(m\) has determinant
      \(1-T^m\); ordinary connectedness does not force \(P_0=1-T\).
\item `Gm`: the localization sequence produces an odd numerator
      \(1-T\) and even denominator \(1-qT\).
\item curves: a \(2g\)-dimensional alternating similitude has determinant
      \(q^g\) and satisfies \eqref{eq:curve-P1-reciprocal}.
\item products: tensor eigenvalues multiply, weights add, and Euler
      characteristics multiply.
\item field extension: Frobenius becomes \(F^m\), and both
      \eqref{eq:base-change-polynomial} and
      \eqref{eq:base-change-zeta} hold.
\end{enumerate}
These tests should depend only on proved cohomology computations and the
algebraic core, and should import no project-specific axioms.
\end{formalization}
